% The Gross-Pitaevskii solver in tnde -- physics user guide.
% Build:  pdflatex guide.tex && pdflatex guide.tex     (twice, for the TOC)
\documentclass[11pt,a4paper]{article}

\usepackage[T1]{fontenc}
\usepackage[utf8]{inputenc}
\usepackage{lmodern}
\usepackage[margin=1in]{geometry}
\usepackage{amsmath,amssymb}
\usepackage{booktabs}
\usepackage{array}
\usepackage{tabularx}
\usepackage{ragged2e}
\usepackage{xcolor}
\usepackage{tcolorbox}
\usepackage{enumitem}
\usepackage{parskip}
\usepackage{listings}
\usepackage[hidelinks]{hyperref}

\definecolor{codebg}{gray}{0.96}
\definecolor{rule}{gray}{0.55}

\lstset{
  basicstyle=\ttfamily\small,
  backgroundcolor=\color{codebg},
  frame=none,
  columns=fullflexible,
  keepspaces=true,
  xleftmargin=1em,
  framexleftmargin=1em,
  framexrightmargin=1em,
  framextopmargin=4pt,
  framexbottommargin=4pt,
  showstringspaces=false,
  breaklines=true,
}

% Inline code.  Underscores are escaped by hand at every call site.
\newcommand{\code}[1]{\texttt{#1}}

% Recurring symbols, so they are typeset identically everywhere.
\newcommand{\kph}{k_{\mathrm{phys}}}
\newcommand{\lFD}{\lambda_{\mathrm{FD}}}
\newcommand{\kcut}{k_{\mathrm{cut}}}
\newcommand{\dd}{\mathrm{d}}

% A boxed statement, for the two findings that must not be skimmed past.
\newtcolorbox{finding}{
  colback=codebg, colframe=rule, boxrule=0.5pt, arc=1pt,
  left=8pt, right=8pt, top=6pt, bottom=6pt,
}

\newcolumntype{L}{>{\RaggedRight\arraybackslash}X}

\title{\textbf{The Gross--Pitaevskii solver in \texttt{tnde}}\\[2pt]
  \large Physics guide: the equation, the splitting and its error budget on quantics tensor trains}
\date{}

\begin{document}
\maketitle
\thispagestyle{empty}

\noindent
What equation is being solved, in what units, with what approximations, and what each
approximation costs. \code{PORTING\_NOTES.md} covers the \emph{engineering} --- index
orders, conventions, benchmarks; this document covers what the numbers mean.

Everything quantitative below was measured against the code in this repository, and
the measurement is named at each claim. The measurements of
Section~\ref{sec:lowpass}, which carry the main finding, are reproduced end to end in
a few seconds by \code{python docs/check\_cutoff.py}, run from the repository root.

\vspace{1em}
\hrule
\vspace{0.5em}
{\small\tableofcontents}
\vspace{0.5em}
\hrule
\newpage

% ---------------------------------------------------------------------------
\section{Running the code}

\begin{lstlisting}[language=Python]
import numpy as np
from tnde import evolve1d, observables

psi0 = lambda x: (1/np.pi)**0.25 * np.exp(-x**2/2)
pots = [lambda x: 0.01*x**2, lambda x: 5.0*np.sin(10*x)**2]

psi, snaps = evolve1d.evolve(psi0, pots, R=30, xmin=-500., xmax=500.,
                             g=5.0, dt=0.01, nsteps=100, maxdim=14)

x, values = observables.reconstruct_window(psi, 30, -500., 500.,
                                           x0=0., halfwidth=20.)
width      = observables.width_dense(psi, 30, -500., 500., halfwidth=50.)
\end{lstlisting}

\noindent
2D is \code{tnde.evolve2d.evolve}, with the same shape of call plus
\code{ymin}/\code{ymax}. Worked examples reproducing the paper's two configurations
are \code{examples/reproduce\_paper\_1d.py} and
\code{examples/reproduce\_paper\_2d.py}; \code{./run\_tests.sh} runs the gates.

Read Section~\ref{sec:lowpass} before trusting a number out of a configuration other
than the paper's two.

% ---------------------------------------------------------------------------
\section{The equation}

The Gross--Pitaevskii equation for a single condensate wave function, in units with
$\hbar = 1$:
\begin{equation}
  i\,\partial_t \psi(x,t)
  = \left[ -\frac{1}{2m}\,\partial_x^2 \;+\; V(x) \;+\; g\,|\psi(x,t)|^2 \right]\psi(x,t),
  \label{eq:gp}
\end{equation}
with $\psi$ normalised to one,
\begin{equation}
  \int |\psi(x,t)|^2\,\dd x = 1 .
\end{equation}
There is no separate particle number: $N$ is absorbed into $g$, which is the only
interaction parameter. $g > 0$ is repulsive.

Nothing in the code carries dimensions. $m$, $g$, $V$ and the box are pure numbers,
and the length and time units are whatever makes $\hbar = 1$ and the given $m$ true.
For the paper's runs $m = 1$, so lengths are in oscillator units of a trap of unit
frequency and times in inverse units of the same.

Every sign convention is fixed in \code{dense.py}, which is the clearest statement of
the scheme in the repository:

\begin{center}
\begin{tabular}{@{}lll@{}}
\toprule
term & propagator over a step $h$ & code \\
\midrule
interaction & $\exp\!\left(-i\,g\,|\psi|^2 h\right)$    & \code{dense.nonlinear\_step} (\code{dense.py:106}) \\
potential   & $\exp\!\left(-i\,V(x)\,h\right)$          & \code{dense.potential\_phase} (\code{dense.py:115}) \\
kinetic     & $\exp\!\left(-i\,k^2 h / 2m\right)$       & \code{dense.kinetic\_factor} (\code{dense.py:75}) \\
\bottomrule
\end{tabular}
\end{center}

The interaction propagator is exact for a time step, not approximate: $|\psi|^2$ is
conserved by the term it generates, so $\exp(-i g |\psi|^2 h)$ is the exact solution
of $i\,\partial_t\psi = g|\psi|^2\psi$ over $h$. The splitting error comes entirely
from the kinetic term not commuting with the two real-space terms.

% ---------------------------------------------------------------------------
\section{The scheme}
\label{sec:scheme}

Second-order Strang splitting, mixed spectral / real space. $K$ is the kinetic
propagator (diagonal in momentum), $V$ the potential phase and $N$ the nonlinear
phase (both diagonal in position), $R$ the renormalisation. Reading \emph{left to
right in time} --- which is the opposite of operator-product convention, so this is
notation, not an operator expression (\code{evolve1d.py:8}):

\begin{lstlisting}
(V/2)(N/2) R  [ K R N V ]^nsteps  K R (N/2)(V/2) R          # 1D
(V/2)(N/2)    [ K R V N ]^nsteps  K R (V/2)(N/2)            # 2D
\end{lstlisting}

$V$ and $N$ are both diagonal in $x$ and commute, so their order --- swapped between
the 1D and 2D drivers --- is cosmetic. The two orders differ only because the Julia
originals differ; they are kept apart deliberately.

Three consequences worth knowing before plotting anything.

\paragraph{A run of \code{nsteps} covers $(\text{\code{nsteps}}+1)\,\dd t$, not
$\text{\code{nsteps}}\cdot \dd t$.}
There are $\text{\code{nsteps}}+1$ kinetic applications: one per loop iteration plus
the closing one. The paper's 1D run, $\text{\code{nsteps}} = 100$ at
$\dd t = 0.01$, evolves to $t = 1.01$.

\paragraph{Snapshot $j$ is not the state at $t = j\,\dd t$.}
It is taken mid-Trotter, at the end of loop iteration $j$. At that point the kinetic
propagator has been applied $j$ times, but $V$ and $N$ have each been applied for
$(j + \tfrac{1}{2})\,\dd t$. So a snapshot is half a step ahead in the real-space
terms and cannot be compared with an exact solution at $t = j\,\dd t$ better than
$O(\dd t)$. This is the convention the reference \code{.jld2} files use, so it is
preserved.

If you need a properly symmetrised state at $t = j\,\dd t$, run \code{evolve} with
$\text{\code{nsteps}} = j-1$ and take the \emph{returned final state}, not a
snapshot. That path applies $K$ a total of $j$ times ($j-1$ in the loop plus the
closing one) and accumulates
\begin{equation*}
  \tfrac{1}{2}\dd t + (j-1)\,\dd t + \tfrac{1}{2}\dd t = j\,\dd t
\end{equation*}
of both $V$ and $N$ --- symmetrised, at $t = j\,\dd t$.

\paragraph{Trotter error is $O(\dd t^2)$ per step and $O(\dd t^2)$ accumulated}
for the symmetrised end-to-end evolution --- but the mid-Trotter snapshots are only
$O(\dd t)$ accurate, for the reason above.

% ---------------------------------------------------------------------------
\section{Discretisation, and the dispersion relation}
\label{sec:dispersion}

The box $[\code{xmin}, \code{xmax}]$ of length $L$ carries $M = 2^R$ points. The 1D
grid includes the endpoint (spacing $L/(M-1)$), the 2D grid does not (spacing $L/M$);
this inconsistency is inherited from the original and documented in
\code{PORTING\_NOTES.md}.

The kinetic step does \textbf{not} use the spectral dispersion $-k^2$. It uses the
eigenvalue of the three-point second difference at spacing $a = L/M$ --- see
\code{exp\_lap\_lowpass\_mpo} in \code{operators.py} and \code{kinetic\_factor}
in \code{dense.py}:
\begin{equation}
  \lFD(k) \;=\; -\frac{4M^2}{L^2}\,\sin^2\!\left(\frac{\pi k}{M}\right),
  \qquad k = 0,\dots,M-1 .
  \label{eq:lamfd}
\end{equation}

Two things follow, and both matter physically.

\paragraph{The DFT index maps to physical momentum as $\kph = 2\pi k / L$.}
Expanding \eqref{eq:lamfd} at small $k$,
\begin{equation}
  \lFD(k) \;\longrightarrow\; -\left(\frac{2\pi k}{L}\right)^{\!2} = -\,\kph^2 .
\end{equation}
Verified numerically for $k = 1,\dots,7$ at $R = 16$, $L = 1000$: the deviation is
$3.8\times 10^{-8}$, which is the expected $O\!\left((\kph a)^2\right)$ residual and
nothing else. This map is independent of $R$ --- it depends only on the box length
--- which is why a cutoff calibrated at one resolution transfers to another.

\paragraph{The dispersion saturates instead of growing.}
$\lFD$ is bounded below by $-4/a^2$ at the Nyquist edge, where the true $-\kph^2$
would be $-(\pi/a)^2 \approx -9.87/a^2$. The error is the standard
$O\!\left((\kph a)^2\right)$ finite-difference one: $|\lFD|$ is $8.1\%$ too small at
$\kph a = 1$ and a factor $2.5$ too small at the Nyquist edge. High-momentum
components therefore propagate too slowly. At the paper's 1D parameters
$a = 9.3\times 10^{-7}$, so this is far below every other error in the calculation
--- but it is the reason the dense oracle and the tensor-train solver agree to
$10^{-12}$ rather than only to the spectral answer, and it must not be ``fixed'' to
$-k^2$ without changing both sides.

% ---------------------------------------------------------------------------
\section{The momentum low-pass --- the dominant approximation}
\label{sec:lowpass}

Wrapped around the kinetic phase is a Fermi--Dirac window in DFT index space,
built by \code{lowpass} in \code{operators.py}. Writing $\sigma$ for the logistic
function $\sigma(z) = 1/(1+e^{-z})$,
\begin{equation}
  W(k) \;=\; 1 \;+\; \sigma\big({-}(k - \kcut)\,\beta\big)
       \;-\; \sigma\big({-}(k - (k_{\max} - \kcut))\,\beta\big) .
  \label{eq:lowpass}
\end{equation}
$W$ is $1$ at both \emph{ends} of the index range --- where the small physical
momenta live, since the DFT wraps negative momenta to high indices --- and $0$ across
the middle. It keeps
\begin{equation}
  |\kph| \;<\; \frac{2\pi\,\kcut}{L}
  \label{eq:physcut}
\end{equation}
and discards the rest.

\textbf{This is not physics.} It is there so the momentum-space operator has a small
bond dimension: a phase $\exp(-i k^2 \dd t / 2m)$ that oscillates all the way to the
Nyquist edge is not compressible, while one that is smoothly switched off is. The
cost is that the kinetic step is no longer unitary --- it removes norm every time it
is applied --- and the \code{normalise} call immediately after every kinetic
application puts that norm straight back.

\begin{finding}
\textbf{Norm conservation is not an accuracy diagnostic in this code.} The norm is
restored by construction at every step. A run that is discarding a percent of the
wave function per hundred steps reports a norm of $1.000000$ throughout.
\end{finding}

\subsection{What the cutoff actually is, and what it costs}

\code{kinetic\_mpo} searches for \code{kcut}: it starts at $2^8$, applies the
operator to a fixed Gaussian trial state $(1/\pi)^{1/4}\exp(-x^2/2)$, and doubles
until the trial retains $99\%$ of its norm. In 2D there is no search ---
$\code{kcut} = 2^8$ is hardcoded, as in the original.

The criterion is calibrated on that trial Gaussian, whose momentum spread is
$\sigma_k = 1$. It knows nothing about the potential the run will actually use. For
the paper's 1D configuration the search therefore stops at:

\begin{center}
\begin{tabular}{@{}lll@{}}
\toprule
\code{kcut} & $\kph$ cutoff & trial norm retained \\
\midrule
$2^8 = 256$          & $1.61$          & $0.9546$ --- rejected \\
$\mathbf{2^9 = 512}$ & $\mathbf{3.22}$ & $\mathbf{0.99996}$ --- \textbf{accepted} \\
\bottomrule
\end{tabular}
\end{center}

\noindent
(measured with \code{dense.kinetic\_factor} at $R = 16$; confirmed against the cached
$R = 30$ operator in \code{opcache/}, which retains $0.999989$ of the TCI'd trial
state.)

Now compare that with the physics of the run. The 1D lattice is $5\sin^2(10x)$, and
\begin{equation}
  \sin^2(10x) = \tfrac{1}{2}\big(1 - \cos 20x\big),
\end{equation}
so the potential drives the wave function at $\kph = 20$ --- \textbf{a factor of $6$
above the cutoff}. The lattice imprints sidebands in real space at every potential
step, and the next kinetic step deletes them.

Measured with the dense oracle at $R = 18$, running the paper's exact 1D
configuration for $100$ steps:

\begin{center}
\begin{tabular}{@{}llll@{}}
\toprule
\code{kcut} & $\kph$ cutoff & max norm lost per step & cumulative (100 steps) \\
\midrule
$2^9$ (what the search picks) & $3.22$  & $3.2\times10^{-4}$ & $\mathbf{3.2\times10^{-2}}$ \\
$2^{12}$                      & $25.7$  & $1.6\times10^{-8}$ & $1.1\times10^{-6}$ \\
$2^{14}$                      & $102.9$ & $0$                & $-5\times10^{-14}$ \\
\bottomrule
\end{tabular}
\end{center}

\noindent
and the effect on the observable the paper reports:

\begin{center}
\begin{tabular}{@{}ll@{}}
\toprule
 & final width $\sqrt{\langle x^2\rangle - \langle x\rangle^2}$ at step 100 \\
\midrule
$\code{kcut} = 2^9$, dense $R = 18$                    & $1.6686$ \\
reference MPS, $R = 30$ (the paper's own data)         & $1.6421$ \\
$\code{kcut} = 2^{14}$ (cutoff converged), dense $R=18$ & $\mathbf{1.4684}$ \\
\bottomrule
\end{tabular}
\end{center}

\noindent
The quantity $1 - \text{overlap}$ between the $2^9$ and $2^{14}$ runs after 100 steps
is $7.4\times 10^{-3}$.

So the tensor-train solver reproduces the reference to $10^{-10}$ --- and both are
far from the same scheme with the aperture opened: $+11.8\%$ in width for the
reference against the converged dense run, $+13.6\%$ comparing the two dense runs
directly (the two percentages differ only because the first pair also differs in $R$
and in truncation).

\begin{finding}
\textbf{The momentum cutoff, not the bond dimension, is the leading approximation in
the 1D configuration.} The $\code{maxdim} = 14$ truncation and the
$\code{tol} = 10^{-10}$ interpolation error are three to eight orders of magnitude
smaller.
\end{finding}

The lattice sideband is visible in the paper's own committed data, suppressed but not
gone: Fourier-transforming the reference state at step 100 over $|x| < 20$ puts
$1.5\times10^{-4}$ of the norm above $\kph = 20$, peaked at $\kph = -20.1$, against
$3.5\times10^{-14}$ in the initial state. The sideband is regenerated in real space
and re-filtered in momentum space every step, reaching a filtered quasi-steady
amplitude rather than its true one.

The 2D configuration is in the same regime, less severely. $L = 200$ and
$\code{kcut} = 2^8$ give a cutoff of $\kph = 8.04$, against an initial momentum kick
of $k = 5$, a Gaussian spread of $1$, and a lattice at $\kph = 2$. Measured with
\code{dense.evolve\_2d}-equivalent stepping at $R = 11$ over 20 steps:
$7.2\times10^{-3}$ norm lost cumulatively at $2^8$, $2.6\times10^{-8}$ at $2^9$, and
$1 - \text{overlap} = 1.3\times10^{-2}$ between them.

\paragraph{If you use this code for new physics, raise \code{kcut} until an
observable stops moving.}
Pass it explicitly, \code{evolve(..., kcut=2**12)}, which also skips the search. The
operator's bond dimension grows slowly with it --- the cached $R = 30$ operator at
$2^9$ has rank $17$ --- so this is affordable.

% ---------------------------------------------------------------------------
\section{Why a tensor train works here}

A quantics representation writes the grid index in binary and gives each bit its own
tensor: site $1$ selects which half of the box, site $2$ which quarter, down to site
$R$ selecting between adjacent grid points. \textbf{A site is a length scale}, and the
bond dimension across a cut is the amount of correlation between scales coarser and
finer than that cut.

That is why this problem compresses. The 1D state is, to good approximation, a smooth
envelope of width ${\sim}1$ multiplied by a lattice-periodic modulation of period
$\pi/10 \approx 0.31$. Those two structures live at well-separated scales and are
nearly factorised, so a cut between them carries little correlation. In a box of
length $1000$ resolved to $2^{30}$ points, the state is captured by $30$ tensors of
rank at most $14$ --- against $10^9$ amplitudes stored densely. Bond dimension is the
physical diagnostic here: it measures how far the state is from a scale-separated
product, and it is what grows when the dynamics genuinely creates structure.

Two consequences of taking that seriously:

\begin{itemize}[leftmargin=1.2em]
\item \textbf{The $2^{30}$ grid is not resolution for its own sake.} It is what lets a
  box wide enough that the wave function never reaches the boundary ($|x| < 500$)
  coexist with a spacing fine enough to resolve the lattice
  ($a = 9.3\times10^{-7}$, i.e.\ $3.4\times10^5$ points per lattice period). Neither
  is negotiable, and densely they are incompatible.
\item \textbf{The 2D train is interleaved}, $x_1 y_1 x_2 y_2 \dots x_R y_R$
  (\code{batcheval.interleave}). A cut after site $2n$ then separates \emph{all}
  structure coarser than $L/2^n$ from all structure finer, in both directions at once
  --- the physically natural cut for an isotropic state. The alternative ordering,
  all $x$ sites then all $y$, would put a cut between the two variables, and its bond
  dimension would be the full $x$--$y$ entanglement of the state, which for the
  rotated lattice components $\sin^2\!\big((x+y)/\sqrt{2}\big)$ is large.
\end{itemize}

The nonlinearity is the one term that cannot be precomputed:
$\exp(-i g |\psi|^2 \dd t)$ depends on the state, so it is re-interpolated from
scratch at every step (\code{evolve1d.apply\_nonlinearity}). That is the dominant cost
per step, and it is why the interaction --- not the trap or the lattice --- is what
makes this expensive.

% ---------------------------------------------------------------------------
\section{The two configurations, as physics}

\subsection{1D --- a breathing quench in a shallow optical lattice}

\code{examples/reproduce\_paper\_1d.py}; $R = 30$, $x \in [-500, 500]$,
$\dd t = 0.01$, $\code{nsteps} = 100$, $g = 5$, $m = 1$, $\code{maxdim} = 14$.
\begin{align}
  \psi_0(x) &= (1/\pi)^{1/4} e^{-x^2/2}, \\
  V(x)      &= 0.01\,x^2 \;+\; 5\sin^2(10x).
\end{align}

\begin{itemize}[leftmargin=1.2em]
\item \textbf{The trap is not the one that made the state.} $\psi_0$ is the ground
  state of a harmonic trap of frequency $1$; the trap it is released into is
  $V = 0.01x^2 = \tfrac{1}{2}m\omega^2 x^2$ with $\omega = \sqrt{0.02} = 0.1414$,
  period $2\pi/\omega = 44.4$. The trap's own ground-state width is
  $\sqrt{1/(2\omega)} = 1.880$ against the initial $\sqrt{1/2} = 0.707$. The state is
  $2.7\times$ too narrow and expands: this is a breathing quench, which is what
  \code{tests/test\_dense\_breathing.py} gates the dense solver against analytically.
\item \textbf{The interaction is strong and reinforces the expansion.} At the peak,
  $g|\psi|^2 = 5/\sqrt{\pi} = 2.82$, comparable to the lattice depth $5$ and twenty
  times the trap frequency. The measured width goes
  $0.707 \rightarrow 1.046 \rightarrow 1.642$ at steps $1$, $50$, $100$.
\item \textbf{The lattice is shallow.} With $V = A\sin^2(k_L x)$ and $k_L = 10$, the
  recoil energy is $E_R = k_L^2/2m = 50$, so the depth $A = 5$ is $0.1\,E_R$ --- far
  too shallow to trap --- and the sidebands it generates at $\kph = \pm 2k_L = \pm 20$
  are exactly what the low-pass of Section~\ref{sec:lowpass} removes.
\item \textbf{The run is short.} $t = 1.01$ in total, $2.3\%$ of a trap period. This
  is early-time expansion dynamics, not a breathing cycle.
\end{itemize}

\subsection{2D --- a moving wave packet in a four-fold quasi-periodic lattice}

\code{examples/reproduce\_paper\_2d.py}; $R = 20$ per direction (a $10^{12}$-point
grid), with $x, y \in [-100, 100]$, $\dd t = 0.01$, $\code{nsteps} = 20$,
$g = 5$, $\code{maxdim} = 50$, and $\code{kcut} = 2^8$ held fixed.
\begin{align}
  \psi_0(x,y) &= (1/\pi)^{1/4}\,e^{-(x^2+y^2)/2}\,e^{5ix}, \\
  V(x,y)      &= 0.001\,(x^2+y^2) \nonumber\\
              &\quad + 10\left[\sin^2 x + \sin^2 y
                 + \sin^2\!\frac{x+y}{\sqrt{2}}
                 + \sin^2\!\frac{y-x}{\sqrt{2}}\right].
\end{align}

\begin{itemize}[leftmargin=1.2em]
\item $e^{5ix}$ is a \textbf{momentum kick}: the packet moves in $+x$ at velocity
  $k/m = 5$, covering ${\sim}1$ length unit over the run's $t = 0.21$ --- a fraction
  of a lattice period.
\item The potential is \textbf{four sinusoidal modulations, two axis-aligned and two
  rotated by $45^\circ$}, with incommensurate periods ($\pi$ and $\pi\sqrt{2}$). This
  is a quasi-periodic potential, not a crystal: it has no unit cell, which is
  precisely the case where the interleaved quantics representation earns its keep, and
  it is the reason the state's rank grows fastest here.
\item Depth $10$ against $E_R = 1/2$ for $k_L = 1$ --- this lattice is $20\,E_R$,
  deep, unlike the 1D one.
\item The trap $0.001(x^2+y^2)$ gives $\omega = 0.0447$, period $141$ --- over the
  20-step run it is essentially flat, present to keep the packet bound rather than to
  shape the dynamics.
\end{itemize}

% ---------------------------------------------------------------------------
\section{Observables}

\code{observables.py} provides position moments and reconstruction only:

\begin{center}
\begin{tabularx}{\textwidth}{@{}llL@{}}
\toprule
quantity & MPO route & direct-summation route \\
\midrule
$\langle x\rangle$, $\langle x^2\rangle$, width
  & \code{width}, \code{expectation\_value} & \code{width\_dense} \\
\addlinespace
2D $\langle x\rangle$, $\langle y\rangle$, $\langle r\rangle$, $\langle r^2\rangle$, widths
  & \code{moments\_2d} & \code{moments\_2d\_dense} \\
\addlinespace
$\psi(x)$, $|\psi(x,y)|^2$
  & --- & \code{reconstruct}, \code{reconstruct\_window}, \code{density\_2d} \\
\bottomrule
\end{tabularx}
\end{center}

\paragraph{Prefer the \code{\_dense} routes.}
A quantics MPO for $x^2$ is built by cross interpolation to a tolerance
\emph{relative to} $\max|x^2|$, which on the paper's box is $2.5\times10^5$. A
tolerance of $10^{-6}$ therefore buys an absolute accuracy of ${\sim}0.25$ against an
expectation value of order $0.5$. This is why the original's committed
\code{widths.txt} begins at $0.8117$ where the analytic answer is
$\sqrt{1/2} = 0.70711$; \code{width\_dense} returns $0.707107$. The \code{\_dense}
routes reconstruct the wave function on a sample and sum directly, and are limited
only by the sampling density --- the state is still never materialised on the full
grid.

\paragraph{Not implemented, and worth knowing.}
There is no energy and no chemical potential. The natural physical invariants of this
system are the Gross--Pitaevskii energy functional
\begin{equation}
  E[\psi] \;=\; \int \left[ \frac{|\partial_x \psi|^2}{2m}
    \;+\; V\,|\psi|^2 \;+\; \frac{g}{2}\,|\psi|^4 \right] \dd x
  \label{eq:energy}
\end{equation}
and the chemical potential
\begin{equation}
  \mu \;=\; E \;+\; \frac{g}{2}\int |\psi|^4 \,\dd x .
\end{equation}
$E$ is conserved by the exact dynamics and would be the sharpest available check on
the splitting --- considerably sharper than the norm, which as
Section~\ref{sec:lowpass} explains is restored by construction and diagnoses nothing.
Both are computable from quantities the code already has (the kinetic term via the
same Fourier MPO, the rest as diagonal expectation values), but neither is written.

% ---------------------------------------------------------------------------
\section{What each test pins, physically}

\begin{center}
\begin{tabularx}{\textwidth}{@{}lL@{}}
\toprule
test & what it establishes \\
\midrule
\code{test\_dense\_breathing.py} &
  The dense oracle reproduces the analytic breathing law
  $\langle x^2\rangle(t) = \langle x^2\rangle_0\cos^2\omega t
   + \langle p^2\rangle_0 \sin^2\!\omega t/(m\omega)^2$
  for a Gaussian in a harmonic trap. Pins the trap, the dispersion and the Trotter
  order against a closed-form solution. \\
\addlinespace
\code{test\_kinetic.py} &
  The tensor-train kinetic propagator matches the dense FFT. Pins the Fourier sign,
  the unitary normalisation, all three site reversals, the Fermi--Dirac window and the
  finite-difference dispersion at once. \\
\addlinespace
\code{test\_fourier.py} & The quantics Fourier MPO is the DFT. \\
\addlinespace
\code{test\_evolve1d.py}, \code{test\_evolve2d.py} &
  Full evolution against the dense oracle, linear and nonlinear, at small $R$. \\
\addlinespace
\code{test\_reference\_r30.py} &
  The $R = 30$ run against the paper's own committed MPS tensors. \\
\addlinespace
\code{test\_tt.py}, \code{test\_fit.py} &
  Tensor-train algebra and the variational MPO application. \\
\bottomrule
\end{tabularx}
\end{center}

Note what this ladder does and does not establish. Every gate compares the tensor
train against a dense solver \emph{using the same discretisation and the same
low-pass}, or against the original's own output. It shows the tensor-train machinery
is faithful to the scheme. It does not, anywhere, check the scheme against the
continuum Gross--Pitaevskii equation \eqref{eq:gp} --- except
\code{test\_dense\_breathing.py}, which does, at $g = 0$ and $\code{kcut} = 2^{12}$.

% ---------------------------------------------------------------------------
\section{Error budget}

For the paper's 1D configuration, largest first:

\begin{center}
\begin{tabularx}{\textwidth}{@{}lLl@{}}
\toprule
source & size & controlled by \\
\midrule
momentum low-pass & $0.20$ absolute in the width, $+13.6\%$ --- measured
  (Section~\ref{sec:lowpass}) & \code{kcut} \\
\addlinespace
Trotter splitting & $O(\dd t^2) \sim 10^{-4}$ --- estimated, not measured & \code{dt} \\
\addlinespace
MPS truncation & \textbf{not measured} --- would need $\code{maxdim} = 14$ against
  $\code{maxdim} \geq 28$ & \code{maxdim}, \code{tolerance} \\
\addlinespace
cross-interpolation & ${\sim}10^{-10}$, but see the caveat below & \code{tolerance} \\
\addlinespace
finite-difference dispersion & $O\!\left((\kph a)^2\right) \sim 10^{-10}$ at $R = 30$
  & \code{R} \\
\bottomrule
\end{tabularx}
\end{center}

The $7.7\times10^{-6}$ max pointwise difference quoted in the README is \textbf{not} a
truncation error: it is the disagreement between this port and the paper's own
tensors, both running at $\code{maxdim} = 14$ with the same truncation, so it bounds
porting fidelity and says nothing about how much the bond-dimension cap costs.

Two traps in reading these:

\begin{itemize}[leftmargin=1.2em]
\item \textbf{TCI's reported error is estimated on its own pivots} and says nothing
  about a region no pivot reached. On the paper's initial state it has been observed
  to report $7\times10^{-11}$ while missing half the wave function. Use
  \code{tt.max\_error} against an independent sample.
\item \textbf{A converged norm means nothing} (Section~\ref{sec:lowpass}).
\end{itemize}

% ---------------------------------------------------------------------------
\section{Symbols}

\begin{center}
\begin{tabularx}{\textwidth}{@{}llL@{}}
\toprule
symbol & code & meaning \\
\midrule
$\psi$   & \code{psi}, \code{psi0} & wave function, $\int|\psi|^2\dd x = 1$ \\
$g$      & \code{g}    & interaction strength; $g > 0$ repulsive \\
$m$      & \code{m}    & mass; $1$ in both paper runs \\
$V$      & \code{potentials} & external potential; a \emph{list}, summed \\
$R$      & \code{R}    & quantics digits; the grid is $2^R$ points per direction \\
$M$      & \code{1 << R} & grid points, $M = 2^R$ \\
$L$      & \code{xmax - xmin} & box length \\
$a$      & $L/M$ or $L/(M-1)$ & grid spacing (Section~\ref{sec:dispersion}) \\
$k$      & index $0\dots M-1$ & DFT index \\
$\kph$   & --- & physical momentum, $2\pi k/L$ (Section~\ref{sec:dispersion}) \\
$\kcut$  & \code{kcut} & low-pass in DFT index; physical cutoff $2\pi \kcut/L$ \\
$\beta$  & \code{beta} & sharpness of the Fermi--Dirac window; $2$ \\
$\dd t$  & \code{dt}   & Trotter step \\
$\chi$   & \code{maxdim} & maximum bond dimension \\
\bottomrule
\end{tabularx}
\end{center}

% ---------------------------------------------------------------------------
\section*{References}
\addcontentsline{toc}{section}{References}

\begin{enumerate}[leftmargin=1.5em]
\item M.~Niedermeier \emph{et al.}, \emph{Solving the Gross--Pitaevskii equation with
  quantics tensor trains}, \href{https://arxiv.org/abs/2507.04262}{arXiv:2507.04262}.
  Original Julia code:
  \href{https://github.com/MarcelNiedermeier/Gross-Pitaevskii-TCI}{MarcelNiedermeier/Gross-Pitaevskii-TCI}.
\item J.~Chen and M.~Lindsey, \emph{Direct interpolative construction of the discrete
  Fourier transform as a matrix product operator},
  \href{https://arxiv.org/abs/2404.03182}{arXiv:2404.03182} --- the construction in
  \code{fourier.py}.
\end{enumerate}

\end{document}
